\documentclass[12pt,hidelinks]{article}
\input{config.tex}

% % --------------------------------------------------------

\begin{document}
\title{Hypothesis test of conditional subsampling}
 \author{HaiYing Wang\\
   Department of Statistics, University of Connecticut}
\maketitle
% \begin{abstract}
% \end{abstract}

\begin{section}{Hypothesis Test}

We note that that the score function of IPW is
\begin{equation*}
    \psi_W (\btheta; \x, y, \delta) = \frac{\delta}{\pi(\x,y)} \nabla_{\btheta} \dl  (y\mid \x,\btheta),
\end{equation*}
and 
\begin{align}\label{eq:score_W}
    \sum_{i=1}^{N} \psi_W (\htheta_W; \x_i, y_i, \delta_i) = 0,
\end{align}
where $\htheta_W$ is the IPW estimator.
Samewise, the score function of MSCLE is 
\begin{align}
    \psi_S (\btheta; \x, y, \delta) &= \delta [\nabla_{\btheta} \dl (y\mid \x,\btheta) - \nabla_{\btheta} \log \bapi(\x, y)],
\end{align}
and
\begin{align}\label{eq:score_S}
    \sum_{i=1}^{N} \psi_S (\htheta_S; \x_i, y_i, \delta_i) = 0,
\end{align} 
 where $\htheta_S$ is the MSCLE and $\bapi(\x, y) = \int f(y \mid \x,\btheta) \pi(\x, y)$.
As it is mentioned by \cite{wang2022maximum}, the IPW estimator $\hbeta_W$ is
consistent, because the conditional expectation of the selection indicator
$\delta$ give the data equals the inclusion probability and therefore the
expectation of the objective function for the IPW estimator is the same as that
of the full data MLE.
Yet, the MSCLE is estimated by the assumed model structure to derive the subsampled conditional likelihood, 
so if the subsampling probabilities also depend on the data then the MSCLE is not consistent to the same limit as the full data MLE when the assumed model structure is misspecified. 
That brings a thought that the distance between IPW estimator and MSCLE can be used as an indicator of model misspecification since this distance vanishes asymptotically in the absence of misspecification, but generally doesn't vanish otherwise.
Hausman-type test \cite{hausman1978specification, white1982maximum} is established to evaluate the difference between these two parameters, i.e., one is stable when no misspecification, while the othe is more efficient yet not that stable.
The misspecification indicator is $\htheta_S - \htheta_W$, so we investigate the asymptotic distribution of $\sqrt{N} (\htheta_S - \htheta_W)$. 
We anticipate that with appropriate regularity conditions, $\sqrt{N} (\htheta_S - \htheta_W)$ will be normally distributed asymptotically with mean zero in the absence of misspecification. 
Given a consistent estimator for the asymptotic covariance matrix, we can form an asymptotic chi-square statistic. 
\begin{assumption}
    \begin{itemize}
        \item[(1)] $|| \psi_{\mathcal{T}} (\btheta; \x, y, \delta)||$ and $||\nabla_{\btheta} \psi_{\mathcal{T}} (\btheta; \x, y, \delta)||$ are integrable, where $\mathcal{T} = \{W,S, SM\}$ and $||A|| = tr^{1/2} (A^\top A)$ is the norm for a vector or matrix $A$.
        \item[(2)] The parameter space $\btheta$ is compact and the third order partial derivative of $\psi_{\mathcal{T}} (y\mid \x,\btheta)$ and $\log \bapi(\x, y)$ with respect to any components of $\btheta$ is bounded in absolute value by an integrable function $B(\x; y)$ that does not depend on $\btheta$.
        \item[(3)]  The variance matrix of the two estimators are finite and positive definite.
        \item[(4)] $\btheta_0$ is the unique solution to the population minimization problem $\min\limits_{\btheta \in \Theta} \E [\dl (y\mid \x,\btheta)]$.
        \item[(5)] Assume that the weighted objection function $\ell_{W} (y\mid \x,\btheta) $ is twice continuously differentiable and is bounded in absolute value by an integrable function $B(\x; y)$ that does not depend on $\btheta$ {(\color{red}Note: This may overlap with Assumption~(2).)}
        \item[(6)] $\Sigma_W^{-1} - \Sigma_S^{-1}$ is positive semidefinite. {\color{red} However, according to the Theorem 2 in \cite{wang2022maximum}, $\Sigma_W^{-1} \geq \Sigma_S^{-1}$, hence this condition is automatically satisfied. }
        
    \end{itemize}
\end{assumption}
Besides, we also make an Assumption 2 that all the properties in Assumption 1 are satisfied for the misspecificated model $f^w$.

Assumptions~(1)--(3) provide the essential regularity conditions for the asymptotic properties of the MSCLE estimator. Assumptions~(3)--(5) furnish the necessary regularity conditions for the asymptotic properties of the IPW estimator. Assumption~(6) ensures the existence of the covariance matrix of $\htheta_S - \htheta_W$, which supports Lemma~\ref{le:test stat}.

\begin{lemma}\label{le:test stat}
    Under the null hypothesis of no misspecification, we consider the test statistic
    \begin{equation*}
        H_N = N (\htheta_S - \htheta_W) \hV_H^{-1} (\htheta_S - \htheta_W)^\top.
    \end{equation*}
    We have $\lN H_N \xrightarrow{d} \chi^2_p$,
    where $p$ is the parameter dimension and $\hV_H$ is the approximated estimation of covariate matrix of $\sqrt{N}(\htheta_S - \htheta_W)$ that $\lN \hV_H = \Sigma_W^{-1} - \Sigma_S^{-1}$. $\Sigma_W$ and $\Sigma_S$ are the inverse of asymptotic matrixs of IPW estimator and MSCLE respectively.
\end{lemma}


\begin{proof}
Since $\htheta_W$ and $\htheta_S$ all have the asymptotic properties, then under the null hypothesis, it follows from \cite{liu2016inverse, wang2022maximum} that
\begin{equation*}
   \lN \sqrt{N} \htheta_W \xrightarrow{d} \N(\btheta_0, \Sigma_W^{-1}),
\end{equation*} 
and 
\begin{equation*}
   \lN \sqrt{N} \htheta_S \xrightarrow{d} \N(\btheta_0, \Sigma_S^{-1}),
\end{equation*}
where $\btheta_0$ is the true parameter.
% While under the situation of misspecification, two estimators also converge but not converge to the same parameter.
% IPW estimator still converges to the MLE of full data.
% Yet MSCLE is optimized through the expectation of prior distribution $\bapi(\x, y)$, which is related with the model assumption. 
% When the model is misspecificated, a wrong $\bapi(\x, y)$ will lead to a biased MSCLE. 

In this case, we derictly use the Hausman's Lemma
Consider the difference estimator $D_N = \htheta_S - \htheta_W$, 
we first have $\lN \sqrt{N} D_N \xrightarrow{d} \N (0, V_H)$, that the approximation estimate of $V_H$ can be derived as
\begin{equation*}
    \hV_H = \Var(\sqrt{N} (\htheta_S - \htheta_W)) \xrightarrow{d} \Sigma_W^{-1} + \Sigma_S^{-1} - 2N \Cov(\htheta_W, \htheta_S).
\end{equation*}
Then we just need to calculate the limit form of $\Cov(\htheta_W, \htheta_S)$.

$\htheta_W$ and $\htheta_S$ are the solutions of  \eqref{eq:score_W} and \eqref{eq:score_S}, so we expense these two equations at the true value $\btheta_0$ and have
\begin{align}\label{eq:tay_psi}
    0 &= \frac{1}{\sqrt{N}} \sum_{i=1}^{N} \psi_W (\btheta_0; \x_i, y_i, \delta_i) + \left[\frac{1}{N} \sum_{i=1}^{N} \nabla_{\btheta} \psi_W (\btheta_0; \x_i, y_i, \delta_i) \right] \sqrt{N} (\htheta_W - \btheta_0) + \op, \\
    0 &= \frac{1}{\sqrt{N}} \sum_{i=1}^{N} \psi_S (\btheta_0; \x_i, y_i, \delta_i) + \left[\frac{1}{N} \sum_{i=1}^{N} \nabla_{\btheta} \psi_S (\btheta_0; \x_i, y_i, \delta_i) \right] \sqrt{N} (\htheta_S - \btheta_0) + \op. 
\end{align}
According to the Assumpution 1(1), we assume that $A_W = \Exp [ \nabla_{\btheta} \psi_W (\btheta_0; \x, y, \delta)]$ and $A_S = \Exp [ \nabla_{\btheta} \psi_S (\btheta_0; \x, y, \delta)]$, with the large law of number, we have 
\begin{equation*}
    \frac{1}{N} \sum_{i=1}^{N} \nabla_{\btheta} \psi_W (\btheta_0; \x_i, y_i, \delta_i) \cvp A_W,
\end{equation*}
and 
\begin{equation*}
    \frac{1}{N} \sum_{i=1}^{N} \nabla_{\btheta} \psi_S (\btheta_0; \x_i, y_i, \delta_i) \cvp A_S.
\end{equation*}
Besides,
\begin{equation*}
    \left[\frac{1}{N} \sum_{i=1}^{N} \nabla_{\btheta} \psi_W (\btheta_0; \x_i, y_i, \delta_i) \right] \left[\frac{1}{N} \sum_{i=1}^{N} \nabla_{\btheta} \psi_S (\btheta_0; \x_i, y_i, \delta_i) \right] \xrightarrow{P} \Exp [\psi_W (\btheta_0; \x, y, \delta) \psi_S (\btheta_0; \x, y, \delta)^\top] = B_{WS},
\end{equation*}
then we have
\begin{align}
    \notag
    B_{WS} &= \Exp \left[ \frac{\delta}{\pi(\x,y)} \nabla_{\btheta} \dl (y\mid \x,\btheta_0) \left( \nabla_{\btheta} \dl (y\mid \x,\btheta_0) - \nabla_{\btheta} \log \bapi(\x, y)  \right)^\top \right] \\ 
    &= \Exp_{(\x,y)} \left[ \Exp  \left[ \frac{\delta}{\pi(\x,y)} \nabla_{\btheta} \dl (y\mid \x,\btheta_0) \left( \nabla_{\btheta} \dl (y\mid \x,\btheta_0) - \nabla_{\btheta} \log \bapi(\x, y)  \right)^\top \right] \mid (\x,y)\right]. 
\end{align}
Since $\Exp[\delta \mid \x,y] = \pi(\x,y)$,
\begin{align}\label{B_WS}
    \notag
    B_{WS} &= \Exp_{(\x,y)} \left[ \nabla_{\btheta} \dl (y\mid \x,\btheta_0) \left( \nabla_{\btheta} \dl (y\mid \x,\btheta_0) - \nabla_{\btheta} \log \bapi(\x, y)  \right)^\top \right] \\ \notag
    &= \Exp_{(\x,y)} \left[ \nabla_{\btheta} \dl (y\mid \x,\btheta_0) \nabla_{\btheta} \dl (y\mid \x,\btheta_0)^\top \right] - \Exp_{(\x,y)} \left[ \nabla_{\btheta} \dl (y\mid \x,\btheta_0) \nabla_{\btheta} \log \bapi(\x, y)^\top \right] \\ \notag
    &= \Exp_{(\x,y)} \left[ \nabla_{\btheta} \dl (y\mid \x,\btheta_0) \nabla_{\btheta} \dl (y\mid \x,\btheta_0)^\top \right] - \Exp_{(y)} \left[ \Exp_{(\x)} \left[ \nabla_{\btheta} \dl (y\mid \x,\btheta_0) \right] \nabla_{\btheta} \log \bapi(\x, y)^\top  \right] \\
    &= \Exp_{(\x,y)} \left[ \nabla_{\btheta} \dl (y\mid \x,\btheta_0) \nabla_{\btheta} \dl (y\mid \x,\btheta_0)^\top \right],
\end{align}
where we use $\Exp_{(\x)} \left[ \nabla_{\btheta} \dl (y\mid \x,\btheta_0) \right]=0$ in this case.

Besides, with the Assumption 1(1) and (5), and Eq.~\ref{B_WS}, we have
\begin{equation*}
    A_W = - \Exp [-\nabla_{\btheta} \psi_W (\btheta_0; \x, y, \delta)] = - \Exp_{(\x,y)} \left[- \nabla_{\btheta}^2 \ell_W (\btheta_0; \x, y, \delta) \right] = B_{WS},
\end{equation*} 
while
\begin{equation*}
    A_S  = \Exp [-\nabla_{\btheta}^2 \ell_S(\btheta_0; \x, y, \delta)] = \Sigma_S,
\end{equation*}
with the conclusion in \cite{wang2022maximum} (A.12) that $-\frac{1}{N} \nabla_{\btheta}^2 \ell_S(\btheta_0; \x, y, \delta) = \Sigma_S + \op$.

Thus we have
\begin{equation*}
    \lN \sqrt{N} (\htheta_W - \btheta_0) = - A_W^{-1} (\frac{1}{\sqrt{N}} \sum_{i=1}^{N} \psi_W (\btheta_0; \x_i, y_i, \delta_i)) +\op,
\end{equation*}
and 
\begin{equation*}
    \lN \sqrt{N} (\htheta_S - \btheta_0) = - A_S^{-1} (\frac{1}{\sqrt{N}} \sum_{i=1}^{N} \psi_S (\btheta_0; \x_i, y_i, \delta_i)) +\op.
\end{equation*}
Then
\begin{equation*}
    \lN \Cov(\sqrt{N} \htheta_W, \sqrt{N} \htheta_S) = \lN E[(\htheta_W -\btheta_0) (\htheta_S -\btheta_0)^\top] = A_W^{-1} B_{WS} A_S^{-1 \top} = \Sigma_S^{-1}.
\end{equation*}
Recall that
\begin{equation*}
    \hV_H \xrightarrow{P} \Sigma_W^{-1} + \Sigma_S^{-1} - 2N \Cov(\htheta_W, \htheta_S) \xrightarrow{P} \Sigma_W^{-1} - \Sigma_S^{-1}.
\end{equation*}
Thus based on the property of Wald-type statistics \cite{van2000asymptotic} and Slutsky's Theorem, the test statistic $H_N$ follows a chi-square distribution $\chi_p^2$.
\end{proof}

As a result, $H_N$ can be used as a test statistic to decide whether the IPW or the MSCLE estimator should be chosen. 
% Let $\btheta_0^*$ denote the full-data maximum likelihood estimator. Note that $\htheta_W \xrightarrow{p} \btheta_0^*$ holds under any specification, whereas $\htheta_S \xrightarrow{p} \btheta_0^*$ may fail if the model is misspecified.
The corresponding hypothesis test is formulated as follows:
\begin{equation*}
    H_0: \E[\htheta_S - \htheta_W]=0 \leftrightarrow H_1:  \E[\htheta_S - \htheta_S]\neq 0,
\end{equation*}
that null hypothesis $H_0$ represents that the model assumption $\dl (y\mid \x,\btheta)$ is correct while the alternative hypothesis $H_1$ represents that the model is misspecified.
We then compute the $p$-value $p_H = 1 - F_{\chi_p^2}(H_N)$ and compare it with the predetermined significance level $\alpha$.  
If $p_H \leq \alpha$, we reject $H_0$ and choose the IPW estimator for more stable estimation; otherwise, we accept $H_0$ and select the MSCLE estimator for greater efficiency.  
Thus, we define a new combined estimator that incorporates both the IPW and MSCLE estimators via the hypothesis-test indicator:
\begin{align}
    \htheta_N &= \I\{H_N \geq c_\alpha\} \htheta_W + \I\{H_N < c_\alpha\} \htheta_S,
\end{align}
where $c_\alpha$ is the upper confidence limit of level $\alpha$ for the chi-squared test and $\I(\cdot)$ is the indicator function.
We now study the property of the new estimator $\htheta_N$, considering two separate cases. 
The situation should be discussed seperately. 
First, if $H_0$ is accepted, when $n \rightarrow \infty$, we have $P(H_N > c_\alpha) \rightarrow \alpha$, then we have probability $\alpha$ to choose IPW estimator and probability $1-\alpha$ to choose MSCLE. 
When the model is correctly specified or both the IPW and MSCLE estimators converge to the same full data estimator $\btheta_l$, $\htheta_N$ is the unbias estimatior since it is the combination of two unbiased estimator
% $\sqrt{N} (\htheta_N - \btheta_0)= \sqrt{N} (\htheta_W - \btheta_0) \I\{H_N \geq c_\alpha\} + \sqrt{N} (\htheta_S - \btheta_0) \I\{H_N < c_\alpha\}$.
\begin{equation*}
    \E[ \htheta_N - \btheta_0 ] = \E[\I\{H_N \geq c_\alpha\} (\htheta_W - \btheta_0)] + \E[\I\{H_N < c_\alpha\} (\htheta_S - \btheta_0)] = 0.
\end{equation*}
% \begin{align}\label{eq:HNmean}
%     \notag
%     & \E[(\htheta_W - \btheta_0) \I\{H_N \geq c_\alpha\}] = \E[\E[(\htheta_W - \btheta_0) \mid \htheta_W-\htheta_S] \I\{H_N \geq c_\alpha\}] \\
%     \notag
%     &= \E[N \Cov(\htheta_W - \btheta_0, \htheta_W-\htheta_S) \Sigma_{S} (\htheta_W-\htheta_S) \I\{H_N \geq c_\alpha\}] \\ 
%     &= (\Sigma_W^{-1}-\Sigma_S^{-1}) \Sigma_{S} \E[(\htheta_W-\htheta_S) \I\{N(\htheta_W-\htheta_S)V_{H}^{-1}(\htheta_W-\htheta_S)^\top \geq c_\alpha\}] = 0.
% \end{align}
Since
\begin{equation*}
   \sqrt{N} (\htheta_W-\htheta_S) \sim \N(0, \Sigma_S^{-1}).
\end{equation*}
and the sample space of $\{(\htheta_W-\htheta_S) \in \R^k: \I\{H_N \geq c_\alpha\}=1\}$ should be a exterior of ellipsoid, eq~\eqref{eq:HNmean} equals to 0.
Thus, we have $\htheta_N$ is an unbiased estiamtor.

When the model is misspecificated, the score function of MSCLE would be a little different from Eq.~\ref{eq:score_S} 
since we have to calculate the expection $\bapi(\x, y) = \int f(y \mid \x,\btheta) \pi(\x, y)$, which depends on the assumption of $\ell$.
Thus, the score function of MSCLE under model misspecification is 
\begin{align}
    \psi_{SM}(\btheta; \x_i, y_i, \delta_i) = \delta_i [ \dot{\ell}(y_i \mid \x_i,\btheta) -  E^w\{ \dot{\ell}(y_i \mid \x_i,\btheta) \mid \x_i, y_i, \delta_i=1\} ],
\end{align}
where $E^w$ is the expectation under the assumed working model.
Since IPW and MSCLE converges to different parameter estimations, we 
assume that IPW and MSCLE converages to $\btheta_l$ and $\btheta_{ls}$ respectively, which means $\btheta_{ls}$ is the solution of the following estimation equation \citep{wang2022maximum}
\begin{align}
    E^t[\psi_{SM}(\btheta; \x, y, \delta)] = 0,
\end{align}
or
    \begin{align}
        \E^t \left\{ \E^t \left[ \dl(\btheta_{ls}; \x, y) \pi(\x,y;\btheta) \right]  -\frac{\E^t \left[ \pi(\x,y;\btheta) \mid \x \right]}{\E^w \left[ \pi(\x,y;\btheta) \mid \x \right]} \E^w \left[ \dl(\btheta_{ls}; \x, y) \pi(\x,y;\btheta) \right] \right\} = 0,
    \end{align}
    where we use $\E^t$ to emphasize that the expectation is taken with respect to the true data
    distribution.
Under the model misspecification, the variance- correlation matrixs of MSCLE should be different.
%Another difference is that $\hV_H$ has the different structure and should be more complicated. 
%In this case, we use the influence function to calculation the $\hV_H$ and calculation of variance matrixs are also derived.

We expand the Taylor expression of the score function of 
MSCLE, resulting in that 
\begin{equation*}
       0 = \frac{1}{\sqrt{N}} \sum_{i=1}^{N} \psi_{SM} (\btheta_{ls}; \x_i, y_i, \delta_i) + \left[\frac{1}{N} \sum_{i=1}^{N} \nabla_{\btheta} \psi_{SM} (\btheta_{ls}; \x_i, y_i, \delta_i) \right] \sqrt{N} (\htheta_S - \btheta_{ls}) + R_N,  
\end{equation*}
where 
\begin{equation*}
    R_N = \frac{1}{2} \sqrt{N} (\htheta_S - \btheta_{ls}) \left[\frac{1}{N} \sum_{i=1}^{N} \nabla^2_{\btheta} \psi_{SM} (\btheta_{ls}; \x_i, y_i, \delta_i) \right] (\htheta_S - \btheta_{ls})^\top + \op.
\end{equation*}
With the Assumption 1(2), we can prove that 
\begin{equation*}
    R_N \leq \frac{1}{N} d ||\htheta_S - \btheta_{ls}|| = \op.
\end{equation*}
Also we have
\begin{align}
    B_{SM} = \Var (\psi_{SM} (\btheta; \x, y, \delta)^ {\otimes 2}) = \E^t [\delta (\dl(y \mid \x,\btheta) -  E^w\{ \dl(y \mid \x,\btheta) \mid \x, y, \delta \})^{\otimes 2} ],
\end{align}
where $A ^{\otimes 2} =AA^\top$.
Then with the Assumption 2(1), let $A_{SM}^{-1} B_{SM} A_{SM}^{-1} = \Sigma_{SM}^{-1}$
\begin{equation*}
    \sqrt{N} (\htheta_S - \btheta_{ls}) \xrightarrow{d} N(0, \Sigma_{SM}^{-1}),
\end{equation*}
while for IPW, we assume that the the asymptotic property is that 
\begin{equation*}
    \sqrt{N} (\htheta_W - \btheta_l) = \xrightarrow{d} N(0, \Sigma_W^{-1}),
\end{equation*}
where $\Sigma_W^{-1} = A_{W}^{-1} B_{W} A_{W}^{-1} = \E^t[\dl ^{\otimes 2} (\btheta;,\x,y)] ^{-1} \E\left\{ \frac{\dl ^{\otimes 2} (\btheta;,\x,y)}{\pi(\x,y)}\right\} \E^t[\dl ^{\otimes 2} (\btheta;,\x,y)]^{-1}$ and
$B_{W} = \E[\frac{\dl^{\otimes 2} (\btheta; \x, y)}{\pi(\x, y)}]$.

Otherwise, if the model is misspecified, then as $n \to \infty$, $P(H_N \geq c_\alpha) \to 1$, and $\htheta_N$ converges asymptotically to the IPW estimator.
We prove the following lemma.
\begin{lemma}\label{le:mis_conv}
    When the model is misspecificated, and the difference between IPW and MSCLE is $O_{P}(1/\sqrt{N})$, then we have $\htheta_N \xrightarrow{p} \htheta_W$.
\end{lemma}

\begin{proof}
We first assume that under the model misspecification, the IPW and MSCLE converge to different parameter value,
\begin{align}
    \notag
    \sqrt{N} (\htheta_W - \btheta_l) & \xrightarrow{d} Z_W \sim N(0, \Sigma_{W}^{-1}), \\
    \sqrt{N} (\htheta_S - \btheta_{ls}) &\xrightarrow{d} Z_S \sim N(0, \Sigma_{SM}^{-1}).
\end{align}
We notice that when the difference $\htheta_W - \htheta_S$ has a order larger than $\frac{1}{\sqrt{N}}$, and 
    % Based on the conclusion in \cite{van2000asymptotic}, IPW $\htheta_W$ is estimated by solving
    % \begin{align}
    %     \g_W(\htheta_W) &= \E^t \left\{ \psi_W (\htheta_W; \x, y, \delta) \right\} =  \E^t \left\{ \dl(\htheta_W; \x, y) \right\} = 0,
    % \end{align}
    % and MSCLE $\htheta_S$ is estimated by solving
    % \begin{align}
    %     \notag
    %     \g_S(\htheta_S) &= \E^t \left\{ \E^t \left[ \dl(\htheta_S; \x, y) \pi(\x,y;\btheta_0) \right]  -\frac{\E^t \left[ \pi(\x,y;\btheta_0) \mid \x \right]}{\E^w \left[ \pi(\x,y;\btheta_0) \mid \x \right]} \E^w \left[ \dl(\htheta_S; \x, y) \pi(\x,y;\btheta_0) \right] \right\} \\
    %     &= 0,
    % \end{align}
    % where we use $\E^t$ to emphasize that the expectation is taken with respect to the true data
    % distribution and $\E^w$ is the expectation under the assumed working model and $\btheta_0$ is the fixed parameter in calculating the sampling probabilities. To estimate the distence between IPW and MSCLE, we use Taylor expansion, deriving that 
    % \begin{align}
    %     \htheta_S - \htheta_W &= \dot g_S(\htheta_W)^{-1} g_S(\htheta_W) + \Op.
    % \end{align}
    % Thus, in most cases, as long as $\htheta_W$ isn't the solution of $g_S(\btheta)$, then we have $\htheta_W \neq \htheta_S$.
    %Under model misspecification, the asymptotic covariance matrix of $\sqrt{N} (\htheta_S - \htheta_W)$
   %differs from the one under the null hypothesis, since the two estimators converge to distinct pseudo-true parameters. 
   the divergence of the Hausman-type statistic and the consistency of the pre-test estimator remain valid provided that the covariance estimator $\Cov(\htheta_W, \htheta_S)$ is consistent and positive definite. Under these conditions, 
   $\lN H_N = N (\htheta_S - \htheta_W) \hV_H^{-1} (\htheta_S - \htheta_W)^\top \xrightarrow{p} \infty$, which easily derives that $\I\{H_N \geq c_\alpha\} \xrightarrow{p} 1$. Thus we have $\htheta_S \I\{H_N < c_\alpha\} \xrightarrow{p} 0$, and we have $\htheta_N \xrightarrow{p} \htheta_W$.
\end{proof}

{\color{red} \cite{leeb2005model} and \cite{leeb2008can} presented a conclusion that it is difficult or nearly impossible to obtain the uniform asymptotic normalityy of $\htheta_N$. No estimator can consistently estimate the unconditional distribution of the model-selected estimator, not even locally.
One reason is that the asymptotic variances of $\htheta_S$ and $\htheta_W$ differ, and their centers slightly shift under local misspecification, the distribution of $\htheta_N$ often becomes multimodal. 
When the test statistic $H_N$ fluctuates around the critical value $c_\alpha$, the realized sample is intermittently assigned to the ``narrow peak'' (associated with $\htheta_S$) or the ``broad peak'' (associated with $\htheta_W$). 
Besides, the selection process is intrinsically correlated with the estimators themselves, as they are derived from the same subsampled data. 
This dependency induces significant skewness and heavy-tailedness in the finite-sample distribution. 
Specifically, the tail probability of $\htheta_N$ can be substantially larger than that of a standard normal distribution.}


\subsection{A model averaging estimator}

The biggest problem of handling $\htheta_N$ is that it is a uncontinues estimator at the point of threshold $H_N = c_\alpha$. Thus, we use the model averageing method instead, by adding a smooth weight into the estimator.
The new estimator 
\begin{equation*}
    \htheta_{NA} = w(H_N)\htheta_W + (1-w(H_N))\htheta_S,
\end{equation*}
where $w(H_N)$ is a weight function related with $H_N$.
We have three kinds options to choose for $w(\cdot)$.
\begin{itemize}
    \item AIC criteria: $w(H_N) = \I\{H_N \geq c_\alpha\}$, which brings $\htheta_{NA} = \btheta_{N}$;
    \item Sigmoid Smoothing: $w(H_N) = \frac{1}{1+\exp\{-(H_N-c_\alpha)/s\}}$, where $s$ is the smoothing parameter.
    \item Adaptive weighting: $w(H_N) = \begin{cases}
 0, & \text{if} \quad H_N \leq c_\alpha, \\
1-\frac{c_\alpha}{H_N}, & \text{if} \quad H_N > c_\alpha.
\end{cases}$
\end{itemize}

\begin{assumption}
    \begin{itemize}
        \item [(1)] 
        $w(\cdot)$ is a measurable function with the value region of $[0,1]$. $w(\cdot)$ should be almost continuous function so that when $D_N \xrightarrow{d} D$, we have $w(D_N) \xrightarrow{d} w(D)$.
%         \item [(2)] there exists an $\epsilon$, so that
% \begin{equation*}
%     \E^t[||\frac{\delta}{\pi} \dot{\ell}^w(\btheta; \x, y)||^{2+\epsilon}]< \infty,
% \end{equation*}
% and 
% \begin{equation*}
%     \E^t[||\delta (\dot{\ell}^w(\btheta; \x, y) - \nabla_{\btheta} \log \bapi(\x, y)) ||^{2+\epsilon}]< \infty.
% \end{equation*}
%     \item [(3)] As $N \rightarrow \infty$, $A_W = \lN \frac{1}{N} \sum_{i=1}^{N} \nabla_{\btheta} \psi_W (\btheta; \x_i, y_i, \delta_i)$  and $A_{SM} = \lN \frac{1}{N} \sum_{i=1}^{N} \nabla_{\btheta} \psi_{SM} (\btheta; \x_i, y_i, \delta_i)$ goes to a positive-definite matrix in probability.
    \end{itemize}
\end{assumption}
For the assumption~(1), since we commonly consider $w(\cdot)$ as sigmoid or huber function, the condition is satisfied naturely.
% Assumption~(2) These two equations present the moment conditions of score function under IPW and MSCLE, respectively.
% Such assumption tends to garuentee that the fluctuations in the estimator at scale $\sqrt{N}$ will not produce a mass that escapes to infinity, so that we can have
% \begin{align}
%     \sup_{N} \E^t[||\sqrt{N} (\htheta_{NA} - \btheta_l)||^{2+\epsilon}] < \infty.
% \end{align}

As we discussed, IPW and MSCLE may converges to different value $\btheta_l$ and $\btheta_{ls}$. 
We first consider the difference between $\btheta_l$ and $\btheta_{ls}$. When $\btheta_l - \btheta_{ls} = h$ where $h$ is a constant, then we have $\lN \htheta_W - \htheta_S = h$.
As a result, $\lN H_N \rightarrow \infty$.
Let $\Delta = \htheta_{NA}-\btheta_l$ and 
\begin{equation*}
    Z_W = \htheta_W-\btheta_l, \qquad Z_S = \htheta_S-\btheta_l,
\end{equation*}
as the Therome 4.1 in \citep{hjort2003frequentist} claimed, we calculate the MSE $R(h)$ to evaluate the variance of the estimator $\htheta_{NA}$ and the MSE is defined as 
\begin{align}\label{eq:Rh}
    \notag
R(h) &= tr(\E(\Delta \Delta^\top)) = tr(\E[(w(H_N) Z_W + (1-w(H_N)) Z_S)^\top (w(H_N) Z_W + (1-w(H_N)) Z_S)]) \\
&= (1-w(H_N))^2 ||h||^2 + \frac{1}{N} tr(w(H_N)^2 \Sigma_W^{-1} + w(H_N)(1-w(H_N)) \Sigma_{SM}^{-1} + tr(C_{WS})) \notag \\
&\xrightarrow{N \rightarrow \infty} (1-w(H_N))^2 ||h||^2,
\end{align}
where $C_{WS}$ can be considered as $
    \Cov(Z_W, Z_S) = \Cov(\htheta_W, \htheta_S) = A_W^{-1} B_{WS} A_{S}^{-1}$.
As long as $\lN w(H_N) = 1$, then we can garuentee that $\lN R(h)=0$, while most of the $w(H_N)$ cases follows this requirement.
Next we prove that $\btheta_{NA}$ is an unbiased estimator through
\begin{align}
   \lN \E(\htheta_{NA} - \btheta_l) = \lN \E[(w(H_N) Z_W + (1-w(H_N)) Z_S)] \rightarrow \lN \E[w(H_N) Z_W] = 0.
\end{align}

We next consider that the different between $\btheta_l$ and $\btheta_{ls}$ decreases as $N$ increases. 
We assume that $||\btheta_l - \btheta_{ls}|| \leq \frac{M}{\sqrt{N}}$, which is bounded by a value with order of $N^{-1/2}$ and $\delta_N = \sqrt{N} (\btheta_l - \btheta_{ls})$ which is a bounded sequence.
As a result, through Bolzano-Weierstrass Theorem, there exists a subsequence $\{N_{k}\}$ such that when $k \rightarrow \infty$, $\delta_{N_k}$ converages to a constant vector $\delta_s$ and $||\delta_s|| \leq M$.
Thus let $D_N = \sqrt{N}(\htheta_S - \htheta_W)$.
 As a result, 
 \begin{align}
   \sqrt{N_k}(\htheta_{W,N_k} - \htheta_{S,N_k}) \xrightarrow{d} D \sim \N(\delta_s, V_H),
 \end{align}
where $V_H = \Sigma_W^{-1} + \Sigma_{SM}^{-1} - C_{WS} - C_{WS}^\top$ is the limitation correlation matrix of $\htheta_S - \htheta_W$.
The Hausman statistic can be formulated as 
\begin{equation*}
    H_N = D_N \hV_H^{-1} D_N^\top \xrightarrow{p} D V_H^{-1} D^\top \xrightarrow{d} \chi^2_{\lambda},
\end{equation*}
while the center value $\lambda  = \delta_s V_H^{-1} \delta_s^\top$ that can be bounded that
\begin{equation*}
    \delta_s V_H^{-1} \delta_s^\top \leq \lambda_{\text{max}}(V_H^{-1}) M^2,
\end{equation*}
where $\lambda_{\text{max}}(V_H^{-1})$ is the maximum eigenvalue of $V_H^{-1}$.

On the other hand, the difference between to esitmator $\htheta_{NA}$ and $\btheta_l$ can be formulated as
\begin{align}
    &\Delta = \sqrt{N}(\htheta_{NA}-\btheta_l) \notag \\
    &= \sqrt{N}[(w(H_N) (\htheta_W - \btheta_l) + (1-w(H_N)) (\htheta_S - \btheta_l))] = Z_W - (1-w(H_N))(Z_W -Z_S + \delta_N),
\end{align}
where 
\begin{equation*}
    Z_W \sim \N(0, \Sigma_W^{-1}), \qquad
    Z_W-Z_S \sim \N(\delta_s, V_H).
\end{equation*}

As long as $w(H_N)$ is bounded, based on the consistency of IPW and MSCLE, $\htheta_{NA}$ also converages to $\theta_0$.
Similar with the conclusion of Lemma~\ref{le:mis_conv}, we just have to study the property when the difference order is $O_{P}(1/\sqrt{N})$.
We involve the concept of frequentist  model average (FMA), which established aunified framework to address the problems of model selection and model averaging. 

To evaluate the efficiency of $R(h)$, we nex want to prove that $R(h)$ is bounded with respect to $h$, so that we can illustrate that the mixture estimator is superior than MSCLE. 
Then we have the following lemma:
\begin{lemma}\label{R_formula}
\begin{align}
        \lim\limits_{N \rightarrow \infty} R &= tr(\Sigma_{W}^{-1}) + ||\delta_s||^2 tr (\E[(1 - w(\chi_{p+2}^2(\lambda)))^2] V_H) + ||\delta_s||^2 \E[(1 - w(\chi_{p+4}^2(\lambda)))^2] \notag \\ 
        &+ ||\delta_s||^2 \E[(1 - w(\chi_{p+2}^2(\lambda)))] tr((C_{WD} + C_{WD}^\top)) \notag \\
        &+ ||\delta_s||^2 \{\E[(1 - w(\chi_{p+4}^2(\lambda)))^2] - E[(1 - w(\chi_{p+2}^2(\lambda)))]\} tr((C_{WD} + C_{WD}^\top) V_H^{-1}).
\end{align}
{\color{red} Should I prove that $R \leq tr(\Sigma_{W}^{-1})$ for all $N$}
\end{lemma}

\begin{proof}
    We first calculate the limit of $w(\chi_{p+2}^2(\lambda))$.
    We already knows that $\frac{\chi^2(\lambda)}{\lambda} \xrightarrow{p} 1$ as $\lambda \rightarrow \infty$ and when $ \delta_s \rightarrow \infty$, we have $\lim\limits_{\delta_s \rightarrow \infty} \lambda = \lim\limits_{\delta_s \rightarrow \infty} \delta_s V_H^{-1}\delta_s^\top \xrightarrow{p} \infty$.
    Plus, since $w(\lambda) \in [0,1]$,
    thus we have $\lim\limits_{\lambda \rightarrow \infty} w(\chi^2(\lambda)) \rightarrow 1$.
    Based on the dominated convergence theorem, we have $\lim\limits_{h \rightarrow \infty} \E[w(\chi_{p+2}^2(\lambda))] \xrightarrow{p} 1$.

    Assume that $D_h = V_H^{-1/2}(D-h)$, then $H_N = ||D_h||^2$ and $Q = V_H^{1/2} \E(w( ||D_h||^2)^2  ||D_h||^2) V_H^{1/2}$, where $D_h$ follows a normal distribution with mean of $V_H^{-1/2}h$ and variance of $\I_p$.
Based on the lemma of non-center chi-square distribution, we have
\begin{equation*}
    \E(w(||D_h||^2)D_h) = V_H^{-1/2} h \E(w(\chi_{p+2}^2(\lambda))),
\end{equation*}
and 
\begin{equation*}
    \E(w(||D_h||^2)||D_h||^2) = \E(w(\chi_{p+2}^2(\lambda))) \I_p + \E(w(\chi_{p+4}^2(\lambda))) V_H^{-1/2} hh^\top V_H^{-1/2},
\end{equation*} 
where $\lambda$ is non-center parameter and $\lambda = h^\top V_H^{-1} h$ in our concept.

 
    We now divide $R(\delta_s)$ into four items. 
    The first item is 
    \begin{align}
        \lim\limits_{N_k \rightarrow \infty} \E(Z_{W, N_k}^2) = \Sigma_{W, N_k}^{-1},
    \end{align}
    which can be proved derictly through the conclusion of IPW.

    The second item is 
    \begin{align}
        \notag
        & \lim\limits_{N_k \rightarrow \infty} \E((1- w(H_{ N_k}))^2 (Z_{S, N_k}-Z_{W, N_k} - \delta_s) (Z_{S, N_k}-Z_{W, N_k} - \delta_s)^\top) \\
        \notag
        & = V_H^{1/2} (\E[(1 - w(\chi_{p+2}^2(\lambda)))^2] \I_p + \E[(1 - w(\chi_{p+4}^2(\lambda)))^2] V_H^{-1/2} \delta_s \delta_s^\top V_H^{-1/2}) V_H^{1/2} 
        \\
        \notag
        &= \E[(1 - w(\chi_{p+2}^2(\lambda)))^2] V_H + \E[(1 - w(\chi_{p+4}^2(\lambda)))^2] ||\delta_s||^2.
    \end{align}

    For the third item, we have to prove that
    \begin{align}
        \notag  
        & \lim\limits_{N_k \rightarrow \infty} \E[(1-w(H_{N_k})) Z_{W, N_k}(Z_{S, N_k} - Z_{W, N_k}- \delta_s)^\top] \\
        &= \E[ (1-w(||D_{\delta_s}||^2)) C_{WD} V_H^{-1} (D_{\delta_s} + \delta_s) D_{\delta_s}^\top ],
    \end{align}
We first consider the joint distribution of $Z_W$ and $D_h$, we have 
\begin{align}
    \begin{pmatrix}
        Z_W \\
        D_h
    \end{pmatrix}
    \sim N\left(\begin{pmatrix}
        0 \\
        -h
    \end{pmatrix},
        \begin{pmatrix}
        \Sigma_W^{-1} & C_{WD} \\
        C_{WD}^\top & V_H
    \end{pmatrix}\right),
\end{align}
where $C_{WD} = \Sigma_W^{-1} - C_{WS}$.
Then based on the conditional property of multinormal distribution, we can turn $Z_W$ to $Z_W = C_{WD} V_H^{-1} (D_h + h) +\epsilon$, where $\epsilon$ is a standard normal variable which is independce with $D$. 
As a result, the third item in Eq.~\ref{eq:rh} can be formulated as 
\begin{align}
    \notag
    &\E(w(H_N)(Z_W(Z_S-Z_W-h)^\top))\\
    &= C_{WD}V_H^{-1} \E[w(||D_h||^2) (D_h + h) D_h^\top] \notag \\
    &= C_{WD}V_H^{-1} \E[w(||D_h||^2)] (||D_h||^2 + h D_h^\top) \notag \\
    &= C_{WD}V_H^{-1} [V_H \E[(1 - w(\chi_{p+2}^2(\lambda)))] \I_p + \E[(1 - w(\chi_{p+4}^2(\lambda)))] hh^\top ]  \notag\\
    &- C_{WD}V_H^{-1} [hh^\top \E((1- w(\chi_{p+2}^2(\lambda))))]  \notag \\
    &= C_{WD} \E[(1 - w(\chi_{p+2}^2(\lambda)))] + C_{WD}V_H^{-1} hh^\top [\E[(1 - w(\chi_{p+4}^2(\lambda)))^2] - E[(1 - w(\chi_{p+2}^2(\lambda)))]].
\end{align}
Thus, we complete the proof.
\end{proof}

Besides, we can also proof that when $|\delta_s||$ goes to infinity, the MSE of $\htheta_{NA}$ still converges to the one of IPW, showing the robustness of such estimator. 
For the second item, since $1 - w(\chi_{p+2}^2(\lambda)) \xrightarrow{p} 0$, then $V_H^{1/2} (\E[(1 - w(\chi_{p+2}^2(\lambda)))^2] \I_p) \rightarrow 0$. 
We also have 
    \begin{align}
        V_H^{1/2} \E[(1 - w(\chi_{p+4}^2(\lambda)))^2] V_H^{-1/2} ||\delta_s||^2 V_H^{-1/2} V_H^{1/2} = \E[(1 - w(\chi_{p+4}^2(\lambda)))^2] ||\delta_s||^2.
    \end{align}
To ensure that this item converges to 0 and recall that $\lambda = \delta_s^\top V_H^{-1} \delta_s^\top \sim O_{P}(||M||^2)$, we just need to make sure that $w(H_N) \sim O_{P}(H_N^{-1})$.
Sigmoid smoothing and adaptive weighting functions all satisfies this requirement.
For the third item, we have 
\begin{equation*}
        \lim\limits_{||\delta_s|| \rightarrow \infty} C_{WD} \E[(1 - w(\chi_{p+2}^2(\lambda)))^2] \rightarrow 0,
\end{equation*}
and 
\begin{equation*}
    hh^\top [\E[(1 - w(\chi_{p+4}^2(\lambda)))^2] - \E[(1 - w(\chi_{p+2}^2(\lambda)))^2]] \rightarrow 0
\end{equation*}
when $\lambda = \delta_s^\top V_H^{-1} \delta_s \sim O_{P}(||\delta_s||^2)$.
Then we prove that the third item also converges to 0, so as to the fourth part.

We assume that $w^*(H_N) = \arg\min_{w} R$.
Thus, we have
\begin{lemma}
    We can obtain the optimal function $w^*$ through 
    \begin{align}
        w^*(H_N) &= 1- \frac{f_{\chi_{p+2}^2(\lambda)}(H_N) tr(C_{WD} + C_{WD}^\top) (tr(V_H^{-1})-1)}{2 \left[f_{\chi_{p+2}^2(\lambda)}(H_N) tr(V_H) + f_{\chi_{p+4}^2(\lambda)}(H_N) tr\{(C_{WD} + C_{WD}^\top) V_H^{-1}\} \right]} .
    \end{align}

\end{lemma}
\begin{proof}
    To minimize $R$, we differentiate $R$ with respect to $w(H_N)$, then we have the following equation. 
    \begin{align}\label{R_differ}
        \notag
        & -2(1-w(H_N)) [f_{\chi_{p+2}^2(\lambda)}(H_N) tr(V_H) + f_{\chi_{p+4}^2(\lambda)}(H_N) tr\{(C_{WD} + C_{WD}^\top) V_H^{-1}\}] \\
        &= f_{\chi_{p+2}^2(\lambda)}(H_N) tr(C_{WD} + C_{WD}^\top) (tr(V_H^{-1})-1).
    \end{align}
    By solving Eq.~\ref{R_differ}, we can obtain the optimized $w*(H_N)$.
\end{proof}


On the other hand, the MSE of MSCLE 
\begin{align}
    R_S(h) = N \E[(\hbeta_S - \btheta_l)(\hbeta_S - \btheta_l)]^\top
\end{align},
while we have
\begin{equation*}
        \sqrt{N} (\hbeta_S - \btheta_l) =  \sqrt{N} (\hbeta_S - \btheta_{ls}) + \sqrt{N} (\htheta_S - \btheta_l).
\end{equation*}
Thus 
\begin{equation*}
    R_S(h) = \E[(Z_s-h)(Z_s-h)^\top] = \E[Z_S Z_S^\top] + hh^\top = V_S+ hh^\top.
\end{equation*}
Then as $||h|| \rightarrow \infty$, $R_S(h) \rightarrow \infty$.

% and 
% \begin{equation*}
%     (\htheta_S - \btheta_{ls}) = -\A_{SM}^{-1} \left( \frac{1}{N} \sum_{i=1}^{N} \psi_{SM} (\btheta_{ls}; \x_i, y_i, \delta_i) +\op \right),
% \end{equation*}
% where the Hessian matrix of IPW $\M_W = \E^t[ \frac{\delta}{\pi(\x,y)} \ddot{\ell}(\btheta_l; \x,y) ]$ and the Hessian matrix of MSCLE $\M_S$ is approximated through 
%  \begin{align}
%     M_{S,N} = \frac{1}{N} \sum_{i=1}^{N} \delta \{ \ddot{\ell}^w(\btheta_l; \x_i,y_i) - \E^w(\ddot{\ell}(\btheta_l; \x_i,y_i) \mid \x_i, \delta=1) - \Var^w(\dl (\btheta_l; \x_i,y_i) \mid \x_i, \delta=1)\} .
% \end{align}
% Under the Assumption 2 and the strong law of large numbers, we have $M_{S,N} \xrightarrow{a.s} \M_S$ where
% \begin{align}
%     \M_S = \E^t [ \delta \{ \ddot{\ell}^w(\btheta_l; \x,y) - \E^w(\ddot{\ell}(\btheta_l; \x,y) \mid \x, \delta=1) - \Var^w(\dl (\btheta_l; \x,y) \mid \x, \delta=1)\} ].
% \end{align}

% We also have to make the assumption that the limiting matrix $\M_S$ is non-singular and positive definite.
% We want to calculate $\hat{V}_H$ and 
% \begin{align}
%     \hV_H \xrightarrow{P} \Sigma_{WM}^{-1} + \Sigma_{SM}^{-1} - \Cov(\htheta_W, \htheta_S) - \Cov(\htheta_S, \htheta_W),
% \end{align}
% where the variance-correlation matrix under the model misspecification is 
% \begin{align}
%     \Sigma_{WM}^{-1} = \E^t[\dl ^{\otimes 2} (\btheta;,\x,y)] ^{-1} \E\left\{ \frac{\dl ^{\otimes 2} (\btheta;,\x,y)}{\pi(\x,y)}\right\} \E^t[\dl ^{\otimes 2} (\btheta;,\x,y)]^{-1}.
% \end{align}

% where 
% \begin{equation*}
%     \dot{\ell}_S (\btheta_{ls};,\x,y) = \delta [\dot{\ell}^w (\btheta_{ls};\x,y) - \E^w\{ \dot{\ell}^w (\btheta_{ls};\x,y) \mid \x, \delta=1 \}].
% \end{equation*}
% Then the $\Cov(\htheta_W-\btheta_l, \htheta_S-\btheta_{ls})$ can be transformed to 
% \begin{align}
%     C_{WS} = \M_W^{-1} \E[ \psi_{W,i} \psi_{S,i}^\top] \M_S^{-\top} &= \M_W^{-1} \E[ \dot{\ell}_W (\btheta_{l};,\x,y) \dot{\ell}_S (\btheta_{ls};,\x,y) ^\top] \M_S^{-\top}
% \end{align}
% where 
% \begin{equation*}
%     \dot{\ell}_W (\btheta_{l};,\x,y) \dot{\ell}_S (\btheta_{ls};,\x,y) ^\top = \dot{\ell}^w (\btheta_{l};\x,y) \cdot [\dot{\ell}^w (\btheta_{ls};\x,y) - \E^w\{ \dot{\ell}^w (\btheta_{ls};\x,y) \mid \x, \delta=1 \}]^\top
% \end{equation*}
% and we denote that $\psi_{W,i} = \psi_W (\btheta_l; \x_i, y_i, \delta_i)$ and $\psi_{S,i} = \psi_S (\btheta_{ls}; \x_i, y_i, \delta_i)$ are the score functions of IPW and MSCLE for simplify.
% The formulation of $\Cov(\htheta_S-\btheta_{ls}, \htheta_W-\btheta_l)$ has the similar result. 

% As a result, we have 
% \begin{align}
%     \notag
%     \sqrt{N} (\htheta_W - \btheta_l) & \xrightarrow{d} Z_W \sim N(0, \Sigma_{WM}^{-1}), \\
%     \sqrt{N} (\htheta_S - \btheta_{ls}) &\xrightarrow{d} Z_S \sim N(0, \Sigma_{SM}^{-1}).
% \end{align}
 
% As a result, 
% \begin{align}
%     \notag
%     R(h) &= \Sigma_{SM}^{-1} +V_H^{1/2} (\E[w(\chi_{p+2}^2(\lambda))^2] \I_p + \E[w(\chi_{p+4}^2(\lambda))^2] V_H^{-1/2} hh^\top V_H^{-1/2}) V_H^{1/2} \\
%     \notag
%     &+ C_{SD}V_H^{-1/2} (\E[w(\chi_{p+2}^2(\lambda))^2] \I_p + \E[w(\chi_{p+4}^2(\lambda))^2] V_H^{-1/2} hh^\top V_H^{-1/2}) V_H^{1/2} \\
%     \notag
%     &- C_{SD}V_H^{-1}(hh^\top \E[w(\chi_{p+2}^2(\lambda))^2]) \\
%     \notag
%     &+ V_H^{1/2} (\E[w(\chi_{p+2}^2(\lambda))^2] \I_p + \E[w(\chi_{p+4}^2(\lambda))^2] V_H^{-1/2} hh^\top V_H^{-1/2}) V_H^{-1/2}C_{SD}^\top \\    
%      &- (hh^\top \E[w(\chi_{p+2}^2(\lambda))^2]) V_H^{-1} C_{SD}^\top.
% \end{align}


% Next, we want to prove that when $||h|| = 0$, $R(0)$ closes to the correlation matrix of MSCLE, which reaches the Cramer-Rao bound of such estimator. 
% \begin{lemma}
%     $\lim_{N \rightarrow \infty} R(0) \rightarrow V_{SM}^{-1}$.
% \end{lemma} 

% \begin{proof}
%     The proof method is similar with the former one.
%     We rewrite $\sqrt{N}(\htheta_{NA}-\btheta_l)$ as 
%     \begin{align}
%         \sqrt{N}(\htheta_{NA}-\btheta_l) = \sqrt{N}[(w(H_N) (\htheta_W - \btheta_l) + (1-w(H_N)) (\htheta_S - \btheta_l))] \notag \\
%     &= Z_S + w(H_N)(Z_W -Z_S +h).
%     \end{align}
%     Based on the previous calculation, we have 
%     \begin{align}
%         R(0) \rightarrow \Sigma_{SM}^{-1} + (V_H + C_{SD} + C_{SD}^\top) \E[w(\chi_{p+2}^2(0))^2],
%     \end{align}
%     and we can also derive that the difference between the correlation matrixs of IPW and MSCLE is
%     \begin{align}
%         \Sigma_{WM}^{-1} \rightarrow \Sigma_{SM}^{-1}  + V_H + C_{SD} + C_{SD}^\top.
%     \end{align}
%     If we choose the appropriate $w(\cdot)$ so that $\E[w(\chi_{p+2}^2(0))^2] \rightarrow 0$, then we can prove that $R(0) \rightarrow \Sigma_{SM}^{-1}$.
%     In this case, $R(0)$ can reach the minimum variance.
%     Besides, we have
%     \begin{equation*}
%         R(0) - \Sigma_{WM}^{-1} \rightarrow -(V_H + C_{SD} + C_{SD}^\top) (1-\E[w(\chi_{p+2}^2(0))^2]),
%     \end{equation*}
%     which means that $\Sigma_{WM}^{-1}$ is great than $R(0)$ in the sense of the Loewner order.

%     When we choose the adaptive weighting smoothing method, $\E[w(\chi_{p+2}^2(0))^2] \rightarrow 0$.
%     If we choose the sigmoid smoothing method and set $c_\alpha$ as a fixed number, then $w(\chi_{p+2}^2(0))$ is a non-zero number. But, if we let $c_\alpha$ is related with $N$ which derive that $\exp(\frac{c_\alpha}{s}) \sim \Op(N^{-1})$, then $w(\chi_{p+2}^2(0))$ also satisfies that $\E[w(\chi_{p+2}^2(0))^2] \rightarrow 0$.
% \end{proof}



% We first have to prove that $N \E[(\htheta_{NA} - \btheta_l)^{\otimes 2}] \rightarrow R(h)$.
% First, under the local misspecificated assumption, we already have 
% \begin{align}
%     \begin{pmatrix}
%         \sqrt{N} (\htheta_W - \btheta_l) \\
%         \sqrt{N} (\htheta_S - \btheta_{ls})
%     \end{pmatrix}
%     \xrightarrow{d}
%     \begin{pmatrix}
%         Z_W \\
%         Z_S
%     \end{pmatrix}
%     \sim N\left(0, \begin{pmatrix}
%         \Sigma_W^{-1} & C_{WS} \\
%         C_{WS} & \Sigma_S^{-1}
%     \end{pmatrix}\right).
% \end{align}
% $\E[w(H_N)^2 ||D_N +h||^2] \xrightarrow{d} \E[w(H)^2 ||D_N +h||^2]$. To make sure this, we have to make sure that $H_N \xrightarrow{d} H$.
% Then $R(h)$ can be simplified as the combination of the variance of $\htheta_S$ and the strategy loss caused by $\htheta_W$ that 
% \begin{align}
%     R(h) = \Sigma_{SM}^{-1} + \E(w(H_N)^2(D_N+h)(D_N+h)^\top) + 
% \end{align}


 %$\htheta_{NA}$ tends under {\color{red} the (2.2) assumption} in distribution to 

% {\color{red} The problems are mainly based on the difference between $\htheta_W$ and $\htheta_S$ instead of the misspecificated function. Yet if we use the misspecificated function, we can't directly derive $\htheta_W$ and $\htheta_S$ and their difference. }

\end{section}


\bibliographystyle{agsm}
\bibliography{ref}


\end{document}
